\begin{tikzpicture}[
  font=\scriptsize,
  tahap/.style={draw, rounded corners, minimum height=1.25cm, text width=1.85cm,
                align=center, font=\scriptsize\bfseries},
  panah/.style={-{Stealth[length=2mm]}, thick}
]
  \node[tahap, fill=blue!8] (t1) at (0,0)
    {{\tiny\mdseries Tahap 1}\\Identifikasi dan susun manfaat};
  \node[tahap, fill=blue!8] (t2) at (2.45,0)
    {{\tiny\mdseries Tahap 2}\\Rencanakan realisasi manfaat};
  \node[tahap, fill=blue!8] (t3) at (4.9,0)
    {{\tiny\mdseries Tahap 3}\\Jalankan rencana manfaat};
  \node[tahap, fill=orange!12] (t4) at (7.35,0)
    {{\tiny\mdseries Tahap 4}\\Tinjau dan evaluasi hasil};
  \node[tahap, fill=green!10] (t5) at (9.8,0)
    {{\tiny\mdseries Tahap 5}\\Temukan manfaat lanjutan};

  \draw[panah] (t1) -- (t2);
  \draw[panah] (t2) -- (t3);
  \draw[panah] (t3) -- (t4);
  \draw[panah] (t4) -- (t5);
  \draw[panah] (t5.south) -- ++(0,-0.6) -| (t1.south);
  \node[font=\tiny, gray, fill=white, inner sep=1pt] at (4.9,-1.23)
    {investasi berikutnya memakai pelajaran dari siklus ini};

  \draw[gray, thick] ([yshift=0.15cm]t1.north west) -- ++(0,0.15)
    -- ([yshift=0.3cm]t2.north east) -- ++(0,-0.15);
  \node[font=\tiny, anchor=south] at ($(t1.north)!0.5!(t2.north)+(0,0.32)$)
    {\textit{business case} disusun di sini};
  \node[font=\tiny, anchor=south] at ([yshift=0.15cm]t3.north)
    {proyek berjalan};
  \node[font=\tiny, anchor=south] at ([yshift=0.15cm]t4.north)
    {pemilik manfaat melapor};
\end{tikzpicture}
